% Pursuing Stacks (cotes 134-2 to 134-8): archive cote 134-7, site page 78 (PDF page 79), typescript page 520.
% typescript. Draft compared: ps6.tex:3886-3944.
% AI-transcribed by Claude Opus 5.5 (claude-opus-5-5, adaptive thinking, effort high), Message Batches API, batch ps4 (msgbatch_012Aoc1NuskATPsyBeKSYK7J), 2026-11-12; not yet checked by eye.
\scanpage{134-7/78}
\mspage{520}

invol\over{x}{v}ed in such extra structure.  It may be noted for instance that,
when taking i=n  hence $\hand{\pi}_i(S^n)=\underline{Z}$ ,  we get an operation of the multiplica-%
tive monoid  $\underline{Z}^{\mathrm{mult}}$ \hand{on} ${}^n\hand{\Pi}_\bullet$,  whose action on the homology groups of ${}^n\hand{\Pi}$,
i.e. on the homotopy groups  $\hand{\pi}_j(S^n_{\xout{\illegible{1}}})$\hand{,} is surely an important extra struc-%
ture on these,  which I hope has been studied by the homotopy people
extensively.  Surely it has been known too for a long time (at any rate
since Artin-Mazur's foundational work on profinite homotopy types) that
when $\hand{\pi}_j(S^n)$ is finite (i.e. practically in all cases except j=n) this
ac\hand{t}ion comes from a continuous action of  the multiplicative \xout{groupoid}
monoid  \xout{$\underline{Z}^\wedge$}
\[ M \ = \ \underline{Z}^{\wedge\mathrm{mult}} \ , \]
where  $\underline{Z}^\wedge$ is the completion of $\underline{Z}$ with respect to the topology defined by
subgroups of finite index:
\[ \underline{Z}^\wedge = \lim \underline{Z}/n\underline{Z} \simeq \hand{\prod}_p \underline{Z}_p \qquad , \]
where in the last member the product is taken over all primes,  and $\underline{Z}_p$
denotes p-adic inte\xout{r}gers.  Taking homotopy types over the ring k = $\underline{Z}^\wedge$,
we should get that this monoid  M operates on the object
\[ {}^n\hand{\Pi}_\bullet \overset{L}{\otimes}_{\underline{Z}}\underline{Z}^\wedge \qquad \text{in} \quad D_\bullet(\underline{\mathbf{Ab}}_{\underline{Z}^\wedge}) \qquad , \]
as this may be interpreted as  $L\Pi_\bullet(S^n_k)$\hand{,} and  $k^{\mathrm{mult}}$  operates on $S^n_k$ =S(n,k\hand{)}
for any ring k\over{.}{,} \hand{(}if indeed  S(n,$\pi$)  depends functorially on the variable
k-module $\pi$\hand{)}.  \deleted{(}At any rate,  it is surely well\unsure{\hand{-}} known that \xout{$\underline{Z}^\wedge$ operat}
M = $\underline{Z}^{\wedge\mathrm{mult}}$  operates on the profinite completions of all spheres.  For
the odd-dimensional sp\over{hh}{he}res,  and taking the restriction of this operati-%
on to the \xout{maximal}\inserted{largest} subgroup
\[ M^\times = Z^{\wedge\times} \simeq \hand{\prod}_p \underline{Z}_p^\times \]
of M,  this group may be viewed as the \inserted[typed]{profinite} Galois group of the maximal
cyclotomic extension of the prime field $\underline{Q}$  (deduced by adjoining\unsure{\hand{,}} all
roots of unity),  and it\hand{'}s action o\over{f}{n}  the profinite completion $S^{n\hand{\wedge}}$ of $S^n$
may be viewed as the canonical Galois action,  when $S^{n=\hand{(}2m\xout{+}-1\hand{)}}$ is inter-%
preted as the homotopy type of affine \xout{m-space}\inserted{(complex)} m-space minus the origin
(which makes a sense as a scheme over \xout{\illegible{1}} the prime field $\underline{Q}$).  This was
the interpretation in my mind,  when stating that the action of $\underline{Z}^{\mathrm{mult}}$
and of it's completion $\underline{Z}^{\wedge\mathrm{mult}}$  on the homotopy groups of spheres was
an important extra structure on these.

\margin{78}
